\documentclass[11pt,oneside]{article}

\pdfminorversion=7
\usepackage{QuizStyle}

\hypersetup{
    pdftitle={20261007 Quiz for MATH200 Differential Equations},
    pdfauthor={Jungwoo Kim},
    pdfsubject={Quiz for MATH200 Differential Equations},
    pdfcreator={LaTeX}
}

\begin{document}

\quiztitle{20261007 Quiz}

\begin{problem}{3.5.4}
Find the general solution of the differential equation
\[
    y''-2y'-3y=-6te^{-t}.
\]
\end{problem}

\begin{solution}
The corresponding homogeneous equation has characteristic equation
\[
    r^2-2r-3=(r-3)(r+1)=0,
\]
so
\[
    y_h=c_1e^{3t}+c_2e^{-t}.
\]
The forcing term is a degree-one polynomial times $e^{-t}$. Since $r=-1$ is a simple characteristic root, the usual trial $(At+B)e^{-t}$ overlaps with the homogeneous solution and must be multiplied by $t$. Thus set
\[
    y_p=t(At+B)e^{-t}.
\]
Substitution gives
\[
    y_p''-2y_p'-3y_p
    =\bigl(-8At+2A-4B\bigr)e^{-t}.
\]
Matching coefficients with $-6te^{-t}$ yields
\[
    -8A=-6,
    \qquad
    2A-4B=0,
\]
so
\[
    A=\frac34,
    \qquad
    B=\frac38.
\]
Therefore
\[
    \boxed{
        y=c_1e^{3t}+c_2e^{-t}
        +\frac38te^{-t}+\frac34t^2e^{-t}
    }.
\]
\end{solution}

\begin{problem}{3.5.5}
Find the general solution of the differential equation
\[
    y''+2y'=5+4\sin(2t).
\]
\end{problem}

\begin{solution}
The characteristic equation of the corresponding homogeneous equation is
\[
    r^2+2r=r(r+2)=0,
\]
so
\[
    y_h=c_1+c_2e^{-2t}.
\]
Because the forcing term is a sum, treat the constant and trigonometric terms as separate subproblems and add the corresponding trial forms. A constant trial would duplicate the homogeneous solution $1$, so the constant forcing requires a factor of $t$. The trigonometric forcing has no such duplication. Hence let
\[
    y_p=At+B\cos(2t)+C\sin(2t).
\]
Then
\[
    y_p''+2y_p'
    =2A+(4C-4B)\cos(2t)+(-4B-4C)\sin(2t).
\]
Matching coefficients with $5+4\sin(2t)$ gives
\[
    2A=5,
    \qquad
    C-B=0,
    \qquad
    -B-C=1.
\]
Thus
\[
    A=\frac52,
    \qquad
    B=C=-\frac12.
\]
Therefore
\[
    \boxed{
        y=c_1+c_2e^{-2t}
        +\frac52t-\frac12\sin(2t)-\frac12\cos(2t)
    }.
\]
\end{solution}

\begin{problem}{3.5.6}
Find the general solution of the differential equation
\[
    y''+2y'+y=4e^{-t}.
\]
\end{problem}

\begin{solution}
The characteristic equation is
\[
    r^2+2r+1=(r+1)^2=0,
\]
so
\[
    y_h=(c_1+c_2t)e^{-t}.
\]
The forcing term is proportional to $e^{-t}$, and $r=-1$ is a characteristic root of multiplicity two. Therefore the trial must be multiplied by $t^2$:
\[
    y_p=At^2e^{-t}.
\]
A direct substitution gives
\[
    y_p''+2y_p'+y_p=2Ae^{-t}.
\]
Hence $2A=4$, so $A=2$. Therefore
\[
    \boxed{
        y=c_1e^{-t}+c_2te^{-t}+2t^2e^{-t}
    }.
\]
\end{solution}

% Boyce Section 3.5, Problem 16(a), treated here as the complete quiz problem.
\Needspace{12\baselineskip}
\begin{problem}{3.5.16}
Determine a suitable form for $Y(t)$ if the method of undetermined coefficients is to be used for
\[
    y''+3y'=2t^4+t^2e^{-3t}+2\sin(3t).
\]
\end{problem}

\begin{solution}
The characteristic equation of the corresponding homogeneous equation is
\[
    r^2+3r=r(r+3)=0,
\]
with roots $0$ and $-3$.

For the polynomial forcing $2t^4$, the usual degree-four polynomial trial overlaps with the homogeneous solution corresponding to $r=0$, so it must be multiplied by $t$. For the term $t^2e^{-3t}$, the exponential factor corresponds to the simple root $r=-3$, so this trial also requires one factor of $t$. The term $2\sin(3t)$ has no duplication with the homogeneous solution and requires both sine and cosine terms. Thus a suitable form is
\[
    \boxed{
    \begin{aligned}
        y_p(t)={}&t(A_0t^4+A_1t^3+A_2t^2+A_3t+A_4)\\
        &+t(B_0t^2+B_1t+B_2)e^{-3t}\\
        &+D\sin(3t)+E\cos(3t).
    \end{aligned}}
\]
\end{solution}

% Boyce Section 3.5, Problem 18(a), treated here as the complete quiz problem.
\begin{problem}{3.5.18}
Determine a suitable form for $Y(t)$ if the method of undetermined coefficients is to be used for
\[
    y''+2y'+2y
    =2e^{-t}+2e^{-t}\cos t+4e^{-t}t^2\sin t.
\]
\end{problem}

\begin{solution}
The characteristic equation of the corresponding homogeneous equation is
\[
    r^2+2r+2=0,
\]
with roots $-1\pm i$. Thus $e^{-t}\cos t$ and $e^{-t}\sin t$ belong to the homogeneous solution family.

The forcing term $2e^{-t}$ alone has no duplication with the homogeneous solution, so it contributes $Ae^{-t}$. The remaining two forcing terms belong to the common family $e^{-t}P_2(t)\cos t$ and $e^{-t}P_2(t)\sin t$. Since differentiation mixes the sine and cosine terms, both degree-two polynomial factors must be included. Because $-1\pm i$ are simple characteristic roots, the entire trigonometric trial must also be multiplied by $t$. Hence a suitable form is
\[
    \boxed{
    \begin{aligned}
        y_p(t)={}&Ae^{-t}\\
        &+t(B_0t^2+B_1t+B_2)e^{-t}\cos t\\
        &+t(D_0t^2+D_1t+D_2)e^{-t}\sin t.
    \end{aligned}}
\]
\end{solution}

% Boyce Section 3.5, Problem 21(a), treated here as the complete quiz problem.
\Needspace{16\baselineskip}
\begin{problem}{3.5.21}
Determine a suitable form for $Y(t)$ if the method of undetermined coefficients is to be used for
\[
    y''-4y'+4y=4t^2+4te^{2t}+t\sin(2t).
\]
\end{problem}

\begin{solution}
The characteristic equation of the corresponding homogeneous equation is
\[
    r^2-4r+4=(r-2)^2=0,
\]
so $r=2$ is a double root.

The polynomial forcing $4t^2$ requires a general quadratic polynomial. For $4te^{2t}$, the usual trial $e^{2t}(B_0t+B_1)$ duplicates the homogeneous solution with multiplicity two, so it must be multiplied by $t^2$. Finally, the forcing $t\sin(2t)$ requires both sine and cosine terms with degree-one polynomial coefficients, and there is no resonance at $\pm2i$. Therefore a suitable form is
\[
    \boxed{
    \begin{aligned}
        y_p(t)={}&A_0t^2+A_1t+A_2\\
        &+t^2(B_0t+B_1)e^{2t}\\
        &+(D_0t+D_1)\sin(2t)\\
        &+(E_0t+E_1)\cos(2t).
    \end{aligned}}
\]
\end{solution}

\begin{problem}{3.6.2}
Use the method of variation of parameters to find a particular solution of
\[
    y''-y'-2y=4e^{-t}.
\]
Then check your answer by using the method of undetermined coefficients.
\end{problem}

\begin{solution}
The corresponding homogeneous equation has characteristic equation
\[
    r^2-r-2=(r+1)(r-2)=0,
\]
so take
\[
    y_1=e^{-t},
    \qquad
    y_2=e^{2t}.
\]
Their Wronskian is
\[
    W(y_1,y_2)(t)
    =\begin{vmatrix}
        e^{-t} & e^{2t}\\
        -e^{-t} & 2e^{2t}
    \end{vmatrix}
    =3e^t\ne0,
\]
so $y_1$ and $y_2$ form a fundamental set of solutions. With
\[
    y_p=u_1y_1+u_2y_2,
    \qquad
    g(t)=4e^{-t},
\]
variation of parameters gives
\[
\begin{aligned}
    u_1'&=-\frac{y_2g}{W(y_1,y_2)}=-\frac43,
    &u_1&=-\frac43t,\\
    u_2'&=\frac{y_1g}{W(y_1,y_2)}=\frac43e^{-3t},
    &u_2&=-\frac49e^{-3t},
\end{aligned}
\]
where zero integration constants have been chosen. Hence
\[
    y_p=-\frac43te^{-t}-\frac49e^{-t}.
\]
The term $-(4/9)e^{-t}$ is homogeneous, so we may choose the simpler particular solution
\[
    \boxed{y_p=-\frac43te^{-t}}.
\]

For the undetermined-coefficients check, $r=-1$ is a simple characteristic root, so use the resonant trial $y_p=Ate^{-t}$. Then
\[
    y_p''-y_p'-2y_p=-3Ae^{-t}=4e^{-t},
\]
which gives $A=-4/3$ and confirms the result.
\end{solution}

\begin{problem}{3.6.11}
Verify that
\[
    y_1(t)=t,
    \qquad
    y_2(t)=te^t
\]
satisfy the corresponding homogeneous equation; then find a particular solution of
\[
    t^2y''-t(t+2)y'+(t+2)y=6t^3,
    \qquad
    t>0.
\]
\end{problem}

\begin{solution}
Direct substitution into the corresponding homogeneous equation gives
\[
\begin{aligned}
    t^2y_1''-t(t+2)y_1'+(t+2)y_1
        &=-t(t+2)+t(t+2)=0,\\
    t^2y_2''-t(t+2)y_2'+(t+2)y_2
        &=t(t+2)e^t\bigl(t-(t+1)+1\bigr)=0.
\end{aligned}
\]
Thus $y_1$ and $y_2$ are solutions of the homogeneous equation.

Before applying variation of parameters, divide by $t^2$; this is valid because $t>0$. The standard form is
\[
    y''-\frac{t+2}{t}y'
    +\frac{t+2}{t^2}y
    =6t,
\]
so the forcing term in the variation-of-parameters formula is $g(t)=6t$. Moreover,
\[
    W(y_1,y_2)(t)
    =\begin{vmatrix}
        t & te^t\\
        1 & (t+1)e^t
    \end{vmatrix}
    =t^2e^t\ne0
    \qquad (t>0),
\]
so $y_1$ and $y_2$ form a fundamental set on this interval. With $y_p=u_1y_1+u_2y_2$,
\[
\begin{aligned}
    u_1'&=-\frac{y_2g}{W(y_1,y_2)}=-6,
    &u_1&=-6t,\\
    u_2'&=\frac{y_1g}{W(y_1,y_2)}=6e^{-t},
    &u_2&=-6e^{-t},
\end{aligned}
\]
where zero integration constants have been chosen. Therefore
\[
    y_p=-6t^2-6t.
\]
Since $-6t$ is a homogeneous solution, being a multiple of $y_1=t$, we may choose
\[
    \boxed{y_p=-6t^2}.
\]
\end{solution}

\Needspace{10\baselineskip}
\begin{problem}{7.1.3}
Transform the differential equation
\[
    u^{(4)}-3u=0
\]
into a system of first-order equations.
\end{problem}

\begin{solution}
Introduce one state variable for each derivative up to order three:
\[
    x_1=u,
    \qquad
    x_2=u',
    \qquad
    x_3=u'',
    \qquad
    x_4=u'''.
\]
This choice makes each derivative of a state variable equal to the next state variable, except for the last equation, which is supplied by the original differential equation. Thus
\[
    \boxed{
    \begin{aligned}
        x_1'&=x_2,\\
        x_2'&=x_3,\\
        x_3'&=x_4,\\
        x_4'&=3x_1.
    \end{aligned}}
\]
\end{solution}

\begin{problem}{7.1.4}
Transform the initial value problem
\[
    u''+2u'+4u=2\cos(3t),
    \qquad
    u(0)=1,
    \qquad
    u'(0)=-2
\]
into an initial value problem for two first-order equations.
\end{problem}

\begin{solution}
Let
\[
    x_1=u,
    \qquad
    x_2=u'.
\]
Then $x_1'=x_2$, while the original differential equation gives
\[
    x_2'=u''=-4u-2u'+2\cos(3t)
    =-4x_1-2x_2+2\cos(3t).
\]
The initial conditions become $x_1(0)=1$ and $x_2(0)=-2$. Hence the required first-order initial value problem is
\[
    \boxed{
    \begin{aligned}
        x_1'&=x_2,\\
        x_2'&=-4x_1-2x_2+2\cos(3t),\\
        x_1(0)&=1,
        \qquad
        x_2(0)=-2.
    \end{aligned}}
\]
\end{solution}

\Needspace{22\baselineskip}
\begin{problem}{7.4.2}
Consider the homogeneous system
\[
    \mathbf{x}'
    =\begin{pmatrix}
        2&-1\\
        3&-2
    \end{pmatrix}\mathbf{x}
\]
and the vector-valued functions
\[
    \mathbf{x}^{(1)}(t)
    =\begin{pmatrix}1\\1\end{pmatrix}e^t,
    \qquad
    \mathbf{x}^{(2)}(t)
    =\begin{pmatrix}1\\3\end{pmatrix}e^{-t}.
\]
\begin{problemsubparts}
    \item Show that the given functions are solutions of the given system of differential equations.
    \item Show that $\mathbf{x}=c_1\mathbf{x}^{(1)}+c_2\mathbf{x}^{(2)}$ is also a solution of the given system for any values of $c_1$ and $c_2$.
    \item Show that the given functions form a fundamental set of solutions of the given system.
    \item Find the solution of the given system that satisfies the initial condition
    \[
        \mathbf{x}(0)=\begin{pmatrix}1\\2\end{pmatrix}.
    \]
    \item Find $W(\mathbf{x}^{(1)},\mathbf{x}^{(2)})(t)$.
    \item Show that the Wronskian $W=W(\mathbf{x}^{(1)},\mathbf{x}^{(2)})$ found in part (e) is a solution of Abel's equation
    \[
        W'=(p_{11}(t)+p_{22}(t))W.
    \]
\end{problemsubparts}
\end{problem}

\begin{solution}
Let
\[
    \mathbf{A}=\begin{pmatrix}2&-1\\3&-2\end{pmatrix}.
\]

\noindent\textup{(a)} For the first vector,
\[
    \mathbf{A}\begin{pmatrix}1\\1\end{pmatrix}
    =\begin{pmatrix}1\\1\end{pmatrix},
\]
so
\[
    \mathbf{A}\mathbf{x}^{(1)}(t)
    =\begin{pmatrix}1\\1\end{pmatrix}e^t
    =(\mathbf{x}^{(1)})'(t).
\]
For the second vector,
\[
    \mathbf{A}\begin{pmatrix}1\\3\end{pmatrix}
    =-\begin{pmatrix}1\\3\end{pmatrix},
\]
and therefore
\[
    \mathbf{A}\mathbf{x}^{(2)}(t)
    =-\begin{pmatrix}1\\3\end{pmatrix}e^{-t}
    =(\mathbf{x}^{(2)})'(t).
\]
Thus both functions are solutions of $\mathbf{x}'=\mathbf{A}\mathbf{x}$.

\medskip
\noindent\textup{(b)} By linearity,
\[
\begin{aligned}
    \mathbf{x}'
    &=c_1(\mathbf{x}^{(1)})'+c_2(\mathbf{x}^{(2)})'\\
    &=c_1\mathbf{A}\mathbf{x}^{(1)}+c_2\mathbf{A}\mathbf{x}^{(2)}\\
    &=\mathbf{A}\bigl(c_1\mathbf{x}^{(1)}+c_2\mathbf{x}^{(2)}\bigr)\\
    &=\mathbf{A}\mathbf{x}.
\end{aligned}
\]
Hence every linear combination $c_1\mathbf{x}^{(1)}+c_2\mathbf{x}^{(2)}$ is also a solution.

\medskip
\noindent\textup{(c)} Their Wronskian is
\[
\begin{aligned}
    W(\mathbf{x}^{(1)},\mathbf{x}^{(2)})(t)
    &=\det\begin{pmatrix}
        e^t&e^{-t}\\
        e^t&3e^{-t}
    \end{pmatrix}\\
    &=3-1=2.
\end{aligned}
\]
Since the Wronskian is nonzero, the two solutions are linearly independent and therefore form a fundamental set of solutions.

\medskip
\noindent\textup{(d)} Since the two functions form a fundamental set, write
\[
    \mathbf{x}(t)=c_1\mathbf{x}^{(1)}(t)+c_2\mathbf{x}^{(2)}(t).
\]
At $t=0$, the initial condition gives
\[
    c_1\begin{pmatrix}1\\1\end{pmatrix}
    +c_2\begin{pmatrix}1\\3\end{pmatrix}
    =\begin{pmatrix}1\\2\end{pmatrix},
\]
so
\[
    c_1+c_2=1,
    \qquad
    c_1+3c_2=2.
\]
Thus
\[
    c_1=c_2=\frac12,
\]
and
\[
    \boxed{
        \mathbf{x}(t)
        =\frac12\begin{pmatrix}1\\1\end{pmatrix}e^t
        +\frac12\begin{pmatrix}1\\3\end{pmatrix}e^{-t}
    }.
\]

\medskip
\noindent\textup{(e)} From part (c),
\[
    \boxed{W(\mathbf{x}^{(1)},\mathbf{x}^{(2)})(t)=2}.
\]

\medskip
\noindent\textup{(f)} Here
\[
    p_{11}(t)=2,
    \qquad
    p_{22}(t)=-2,
\]
so $p_{11}(t)+p_{22}(t)=0$. Since $W(t)=2$ is constant,
\[
    W'(t)=0=(p_{11}(t)+p_{22}(t))W(t).
\]
Thus the Wronskian satisfies Abel's equation.
\end{solution}

\end{document}
